\documentclass[12pt,a4paper]{article}

\usepackage[T1]{fontenc}
\usepackage[utf8]{inputenc}
\usepackage[polish]{babel}
\usepackage{geometry}
\geometry{margin=2.5cm}

\usepackage{amsmath, amssymb, amsfonts}
\usepackage{listings}
\usepackage{xcolor}
\usepackage{graphicx}
\usepackage{float}
\usepackage{booktabs}
\usepackage{hyperref}
\usepackage{enumitem}
\usepackage{parskip}
\usepackage{array}

\usepackage{tikz}
\usetikzlibrary{shapes.geometric, arrows.meta, positioning, fit, calc, backgrounds}

\lstdefinestyle{python}{
  language=Python,
  basicstyle=\ttfamily\small,
  keywordstyle=\color{blue!70!black}\bfseries,
  commentstyle=\color{gray},
  stringstyle=\color{orange!80!black},
  showstringspaces=false,
  breaklines=true,
  numbers=left,
  numberstyle=\tiny\color{gray},
  frame=single,
  rulecolor=\color{gray!40},
  backgroundcolor=\color{gray!5},
  tabsize=4,
}

\tikzstyle{startstop} = [rectangle, rounded corners=8pt, minimum width=3.3cm,
                         minimum height=0.9cm, text centered, draw=black,
                         fill=teal!20, font=\small\bfseries]
\tikzstyle{process}   = [rectangle, minimum width=3.3cm, minimum height=0.9cm,
                         text centered, draw=black, fill=blue!10, font=\small]
\tikzstyle{decision}  = [diamond, aspect=2.2, minimum width=3.3cm,
                         minimum height=0.9cm, text centered, draw=black,
                         fill=orange!20, font=\small]
\tikzstyle{io}        = [trapezium, trapezium left angle=70, trapezium right angle=110,
                         minimum width=3.3cm, minimum height=0.9cm, text centered,
                         draw=black, fill=green!10, font=\small]
\tikzstyle{arrow}     = [thick, ->, >=Stealth]
\tikzstyle{darrow}    = [thick, ->, >=Stealth, dashed]

\begin{document}

\begin{titlepage}
  \centering
  \vspace*{3cm}
  {\LARGE\bfseries pyDpVision \par}
  \vspace{0.5cm}
  {\large Dokumentacja modułu \par}
  \vspace{0.3cm}
  {\huge\ttfamily dpVision/xrayProjection.py \par}
  \vspace{1cm}
  \rule{0.7\textwidth}{0.4pt}
  \vspace{0.5cm}
  \par
  {\large Projekcja RTG w przestrzeni świata dla wielu wolumenów\\
  z rozdzieleniem geometrii, fizyki i prezentacji \par}
  \vfill
  {\small \today}
\end{titlepage}

\tableofcontents
\newpage

\section{Cel modułu}

Moduł \texttt{dpVision/xrayProjection.py} został wydzielony po to, aby:
\begin{itemize}
  \item oddzielić obliczenia projekcji RTG od klasy \texttt{Volumetric},
  \item umożliwić pracę na wielu obiektach sceny jednocześnie,
  \item liczyć projekcję w przestrzeni świata z uwzględnieniem transformacji hierarchicznych,
  \item pozostawić otwartą drogę do późniejszego rozszerzenia o nowe źródła próbek
        (np. siatki trójkątów, źródła hybrydowe, modele materiałowe),
  \item rozdzielić \emph{akwizycję} od \emph{prezentacji}.
\end{itemize}

Na obecnym etapie moduł wspiera przede wszystkim:
\begin{itemize}
  \item źródła wolumetryczne przez adapter \texttt{VolumetricXRaySource},
  \item prosty model tłumienia oparty o wartości CT,
  \item projekcję na płaski detektor,
  \item kilka modeli prezentacji końcowego obrazu,
  \item eksport wyniku do \texttt{PNG}, \texttt{TIFF} i \texttt{DICOM Secondary Capture}.
\end{itemize}

\section{Architektura}

\subsection{Warstwy odpowiedzialności}

Rdzeń modułu został rozdzielony na cztery warstwy:
\begin{enumerate}[label=\arabic*.]
  \item \textbf{Geometria projekcji} -- położenie źródła, orientacja detektora,
        rozdzielczość siatki detektora i krok całkowania.
  \item \textbf{Źródła próbkowania} -- obiekty sceny zdolne odpowiadać na pytanie:
        ``jaki współczynnik osłabienia obowiązuje w zadanym punkcie świata?''.
  \item \textbf{Model fizyczny} -- mapowanie z wartości skalarnej do współczynnika
        osłabienia oraz przekształcenie całki po promieniu do surowego wyniku detektora.
  \item \textbf{Model prezentacji} -- budowa obrazu końcowego do oglądania lub eksportu.
\end{enumerate}

\subsection{Najważniejsze klasy}

\begin{table}[H]
\centering
\caption{Główne klasy modułu \texttt{xrayProjection.py}}
\begin{tabular}{>{\ttfamily}p{4.2cm} p{10cm}}
\toprule
\textbf{Klasa / funkcja} & \textbf{Rola} \\
\midrule
XRayProjectionGeometry & opis źródła i detektora w wybranej przestrzeni odniesienia \\
XRayPhysicsModel & mapowanie \textit{scalar $\rightarrow \mu$} oraz \textit{integral $\rightarrow$ detector value} \\
XRaySampleSource & abstrakcyjne źródło próbek dla projektora \\
VolumetricXRaySource & adapter klasy \texttt{Volumetric} do świata RTG \\
XRayProjector & właściwy projektor promieniowania \\
XRayPresentationModel & abstrakcyjna warstwa prezentacji \\
RawPresentationModel & brak modyfikacji obrazu surowego \\
FilmLikePresentationModel & prosty model ``kliszy'' z odwróceniem i nieliniowością \\
DigitalRadiographyPresentationModel & prosty model cyfrowej prezentacji RTG \\
XRayProjectionQualityProfile & preset jakosci zmieniajacy krok ray-marchingu i probkowanie detektora \\
XRayProjectionConfig & pakiet konfiguracji: geometria, fizyka, prezentacja i profil jakosci \\
XRayScene & kontener zrodel sceny i wygodny punkt wejscia do projekcji \\
XRayProjectionStats & podstawowe statystyki czasu i liczby probek dla jednej projekcji \\
\bottomrule
\end{tabular}
\end{table}

\begin{figure}[H]
\centering
\begin{tikzpicture}[node distance=0.9cm and 1.7cm]
  \node (geo)   [io]       {Geometria projekcji};
  \node (src)   [io, below=of geo] {Źródła próbek};
  \node (phys)  [process, right=of geo] {Model fizyczny};
  \node (proj)  [process, right=of src] {XRayProjector};
  \node (raw)   [process, right=of proj] {Wynik surowy};
  \node (pres)  [process, right=of raw] {Model prezentacji};
  \node (out)   [startstop, right=of pres] {Obraz końcowy};

  \draw [arrow] (geo) -- (proj);
  \draw [arrow] (src) -- (proj);
  \draw [arrow] (phys) -- (proj);
  \draw [arrow] (proj) -- (raw);
  \draw [arrow] (raw) -- (pres);
  \draw [arrow] (pres) -- (out);

  \begin{scope}[on background layer]
    \node[draw=black, rounded corners=8pt, fit=(geo)(src), inner sep=0.35cm, fill=green!4] {};
    \node[draw=black, rounded corners=8pt, fit=(phys)(proj)(raw), inner sep=0.35cm, fill=blue!4] {};
    \node[draw=black, rounded corners=8pt, fit=(pres)(out), inner sep=0.35cm, fill=orange!4] {};
  \end{scope}
\end{tikzpicture}
\caption{Podział na geometrię, akwizycję i prezentację}
\end{figure}

\section{Przestrzenie i transformacje}

\subsection{Idea podstawowa}

Każdy wolumen przechowuje dane w swojej lokalnej siatce wokseli. Projektor pracuje
w przestrzeni świata. Dla każdego punktu promienia wykonywany jest łańcuch odwzorowań:
\[
  \mathbf{p}_{\mathrm{ref}}
  \xrightarrow{T_{\mathrm{world}\leftarrow \mathrm{ref}}}
  \mathbf{p}_{\mathrm{world}}
  \xrightarrow{T_{\mathrm{local}\leftarrow \mathrm{world}}}
  \mathbf{p}_{\mathrm{local}}
  \xrightarrow{\texttt{world\_to\_voxel}}
  \mathbf{p}_{\mathrm{voxel}}.
\]

Interpretacja:
\begin{itemize}
  \item \(\mathbf{p}_{\mathrm{ref}}\) -- punkt opisany w przestrzeni odniesienia geometrii,
  \item \(\mathbf{p}_{\mathrm{world}}\) -- ten sam punkt w globalnym układzie sceny,
  \item \(\mathbf{p}_{\mathrm{local}}\) -- punkt w lokalnym układzie konkretnego źródła,
  \item \(\mathbf{p}_{\mathrm{voxel}}\) -- współrzędne ułamkowe w siatce wokseli.
\end{itemize}

\subsection{Budowa geometrii wysokiego poziomu}

Niskopoziomowe API pozwala nadal definiowac detektor przez:
\begin{itemize}
  \item \texttt{detector\_origin\_ref},
  \item \texttt{detector\_u\_ref},
  \item \texttt{detector\_v\_ref},
  \item \texttt{detector\_shape\_hw}.
\end{itemize}

Jednoczesnie modul udostepnia konstruktor
\texttt{XRayProjectionGeometry.from\_detector\_pose(...)}, ktory przyjmuje:
\begin{itemize}
  \item srodek detektora,
  \item normalna detektora,
  \item wektor ``up'',
  \item rozdzielczosc detektora,
  \item oraz albo rozmiar fizyczny detektora, albo wielkosc piksela.
\end{itemize}

To jest preferowany punkt wejscia dla przyszlego GUI, bo pozwala opisywac
scene parametrami blizszymi rzeczywistej geometrii RTG niz surowe wektory
\(\mathbf{u}_{\mathrm{det}}\) i \(\mathbf{v}_{\mathrm{det}}\).

\subsection{Macierz świata dla wolumenu}

Wolumen może być potomkiem dowolnej hierarchii transformacji. Globalne położenie
źródła wolumetrycznego opisuje macierz:
\[
  T_{\mathrm{world}\leftarrow \mathrm{local}} \in \mathbb{R}^{4 \times 4}.
\]

W adapterze \texttt{VolumetricXRaySource} używana jest jej odwrotność:
\[
  T_{\mathrm{local}\leftarrow \mathrm{world}}
  =
  T_{\mathrm{world}\leftarrow \mathrm{local}}^{-1}.
\]

Dla punktu w świecie:
\[
  \tilde{\mathbf{p}}_{\mathrm{local}}
  =
  T_{\mathrm{local}\leftarrow \mathrm{world}}
  \tilde{\mathbf{p}}_{\mathrm{world}},
\]
gdzie \(\tilde{\mathbf{p}}\) oznacza współrzędne jednorodne.

\begin{figure}[H]
\centering
\begin{tikzpicture}[scale=1.0]
  \draw[->, thick] (0,0) -- (3.8,0) node[below] {$X_{\mathrm{world}}$};
  \draw[->, thick] (0,0) -- (0,3.2) node[left] {$Y_{\mathrm{world}}$};

  \draw[->, thick, blue] (1.2,0.8) -- ++(2.0,0.7) node[above] {$x_{\mathrm{local}}$};
  \draw[->, thick, blue] (1.2,0.8) -- ++(-0.4,1.7) node[left] {$y_{\mathrm{local}}$};
  \filldraw[blue] (1.2,0.8) circle (1.5pt) node[below left] {$O_{\mathrm{local}}$};

  \draw[dashed, gray] (1.2,0.8) -- (1.2,0);
  \draw[dashed, gray] (1.2,0.8) -- (0,0.8);
  \node at (2.5,2.3) {$T_{\mathrm{world}\leftarrow \mathrm{local}}$};
\end{tikzpicture}
\caption{Lokalny układ wolumenu osadzony w przestrzeni świata}
\end{figure}

\section{Geometria projekcji}

\subsection{Detektor}

Detektor opisuje punkt początkowy i dwa wektory bazowe:
\[
  \mathbf{o}_D, \qquad \mathbf{u}_D, \qquad \mathbf{v}_D.
\]

Dla piksela \((r, c)\) środek próbki na detektorze wynosi:
\[
  \mathbf{p}_{D}(r,c)
  =
  \mathbf{o}_D + c\,\mathbf{u}_D + r\,\mathbf{v}_D.
\]

W implementacji wartości te są najpierw zadane w przestrzeni referencyjnej,
a następnie mapowane do świata przez \texttt{reference\_transform}.

\subsection{Źródło promieniowania}

Moduł dopuszcza dwa warianty:
\begin{enumerate}[label=\alph*)]
  \item \textbf{projekcja stożkowa} -- z pozycją źródła \(\mathbf{s}\),
  \item \textbf{projekcja równoległa} -- ze stałym kierunkiem \(\mathbf{d}\).
\end{enumerate}

W wariancie stożkowym kierunek promienia dla piksela \((r,c)\) jest równy:
\[
  \hat{\mathbf{d}}(r,c)
  =
  \frac{\mathbf{p}_{D}(r,c) - \mathbf{s}}
       {\|\mathbf{p}_{D}(r,c) - \mathbf{s}\|}.
\]

\begin{figure}[H]
\centering
\begin{tikzpicture}[scale=0.95]
  \coordinate (S) at (-2.5,0);
  \coordinate (D0) at (3.0,-2.0);
  \coordinate (D1) at (3.0, 2.0);
  \coordinate (P) at (3.0, 0.7);

  \draw[thick] (D0) -- (D1);
  \filldraw (S) circle (1.8pt) node[left] {$\mathbf{s}$};
  \filldraw (P) circle (1.5pt) node[right] {$\mathbf{p}_{D}(r,c)$};

  \draw[arrow] (S) -- (P);
  \draw[dashed] (S) -- (D0);
  \draw[dashed] (S) -- (D1);

  \node at (3.5, 0) {detektor};
  \node at (0.2, 1.2) {$\hat{\mathbf{d}}(r,c)$};
\end{tikzpicture}
\caption{Geometria projekcji stożkowej}
\end{figure}

\section{Model akwizycji}

\subsection{Mapowanie z intensywności do współczynnika osłabienia}

W bieżącej implementacji stosowane jest bardzo proste przybliżenie CT-like.
Nie jest to jeszcze model kliniczny, ale daje poprawną strukturę architektoniczną.

Najpierw z wartości skalarnej \(I_{\mathrm{CT}}\) budowana jest względna gęstość:
\[
  \rho_{\mathrm{rel}}
  =
  \max\!\left(0,\,
    1 + \frac{I_{\mathrm{CT}}}{|HU_{\mathrm{air}}|}
  \right).
\]

Następnie liniowy współczynnik osłabienia:
\[
  \mu
  =
  \left(
    \mu_{\mathrm{air}} + \mu_{\mathrm{water}} \rho_{\mathrm{rel}}
  \right)
  \cdot s_{\mu},
\]
gdzie:
\begin{itemize}
  \item \(\mu_{\mathrm{air}}\) -- wkład powietrza,
  \item \(\mu_{\mathrm{water}}\) -- parametr referencyjny,
  \item \(s_{\mu}\) -- skala osłabienia.
\end{itemize}

\subsection{Całkowanie po promieniu}

Dla jednego promienia próbkujemy punkty:
\[
  \mathbf{p}(t_k) = \mathbf{o}_R + t_k \hat{\mathbf{d}},
  \qquad
  k = 0, \dots, N-1.
\]

Jeżeli scena zawiera wiele źródeł, ich wkłady są sumowane:
\[
  \mu_{\mathrm{tot}}(t_k)
  =
  \sum_{m=1}^{M} \mu_m(t_k).
\]

Całka po promieniu jest aproksymowana numerycznie:
\[
  A
  \approx
  \sum_{k=0}^{N-1} \mu_{\mathrm{tot}}(t_k)\,\Delta s.
\]

W bieżącym \texttt{XRayPhysicsModel} wynik surowy można zwrócić jako:
\begin{enumerate}[label=\alph*)]
  \item \textbf{integral}:
  \[
    I_{\mathrm{raw}} = A,
  \]
  \item \textbf{intensity}:
  \[
    I_{\mathrm{raw}} = \exp(-A).
  \]
\end{enumerate}

\section{Model prezentacji}

\subsection{Po co osobna warstwa prezentacji}

Projektor nie powinien decydować, czy kość ma być jasna, czy ciemna.
Projektor powinien dostarczyć wynik surowy. Dopiero kolejny model decyduje,
jak ten wynik pokazać człowiekowi.

Stąd rozdzielenie:
\[
  \text{akwizycja}
  \;\longrightarrow\;
  I_{\mathrm{raw}}
  \;\longrightarrow\;
  \text{prezentacja}
  \;\longrightarrow\;
  I_{\mathrm{display}}.
\]

\subsection{Dostępne modele}

\paragraph{\texttt{RawPresentationModel}}
Zwraca wynik bez zmian:
\[
  I_{\mathrm{display}} = I_{\mathrm{raw}}.
\]

\paragraph{\texttt{FilmLikePresentationModel}}
Stosuje:
\begin{itemize}
  \item normalizację do zakresu roboczego,
  \item opcjonalne odwrócenie jasności,
  \item prostą regulację kontrastu,
  \item nieliniowość \(\gamma\).
\end{itemize}

\paragraph{\texttt{DigitalRadiographyPresentationModel}}
Stosuje podobny pipeline, ale semantycznie odpowiada cyfrowej prezentacji RTG:
\begin{itemize}
  \item windowing,
  \item invert,
  \item kontrast,
  \item gamma.
\end{itemize}

Najprostsza postać takiego modelu jest:
\[
  I_N
  =
  \mathrm{clip}\!\left(
    \frac{I_{\mathrm{raw}} - I_{\min}}
         {I_{\max} - I_{\min}},
    0, 1
  \right),
\]
\[
  I_C = \mathrm{clip}\!\left(0.5 + (I_N - 0.5) c,\; 0,1\right),
\]
\[
  I_{\mathrm{display}}
  =
  \begin{cases}
    (1 - I_C)^{1/\gamma}, & \text{gdy aktywne invert},\\
    I_C^{1/\gamma}, & \text{w przeciwnym razie.}
  \end{cases}
\]

\section{Eksport wyników}

Moduł wspiera trzy typy zapisu:
\begin{enumerate}[label=\arabic*.]
  \item \textbf{PNG} -- szybki podgląd 8-bitowy.
  \item \textbf{TIFF} -- zapis precyzyjny:
    \texttt{uint8}, \texttt{uint16} albo \texttt{float32}.
  \item \textbf{DICOM Secondary Capture} -- prosty eksport do workflow medycznego.
\end{enumerate}

Ważne rozróżnienie:
\begin{itemize}
  \item \texttt{float32 TIFF} jest nośnikiem danych surowych lub półsurowych,
  \item \texttt{PNG}, \texttt{uint16 TIFF} i \texttt{DICOM} zwykle zapisują obraz już
        przekształcony przez model prezentacji.
\end{itemize}

\section{Schemat algorytmu}

\begin{figure}[H]
\centering
\begin{tikzpicture}[node distance=0.75cm and 1.8cm]
  \node (start) [startstop] {Start};
  \node (geo)   [process, below=of start] {Walidacja geometrii};
  \node (pix)   [process, below=of geo] {Wybór piksela detektora};
  \node (ray)   [process, below=of pix] {Budowa promienia w świecie};
  \node (box)   [decision, below=of ray, yshift=-0.1cm] {Przecięcie z AABB sceny?};
  \node (samp)  [process, below=of box, yshift=-0.2cm] {Próbkowanie wszystkich źródeł};
  \node (intg)  [process, below=of samp] {Całkowanie osłabienia};
  \node (raw)   [process, below=of intg] {Wynik surowy piksela};
  \node (next)  [decision, below=of raw, yshift=-0.1cm] {Kolejny piksel?};
  \node (pres)  [process, right=4.2cm of raw] {Model prezentacji};
  \node (save)  [process, below=of pres] {Eksport / wizualizacja};
  \node (stop)  [startstop, below=of save] {Koniec};
  \node (skip)  [process, right=4.2cm of box] {Piksel pusty};

  \draw [arrow] (start) -- (geo);
  \draw [arrow] (geo) -- (pix);
  \draw [arrow] (pix) -- (ray);
  \draw [arrow] (ray) -- (box);
  \draw [arrow] (box) -- node[left] {tak} (samp);
  \draw [arrow] (box) -- node[above] {nie} (skip);
  \draw [arrow] (skip) |- (next);
  \draw [arrow] (samp) -- (intg);
  \draw [arrow] (intg) -- (raw);
  \draw [arrow] (raw) -- (next);
  \draw [arrow] (next) -- node[right] {tak} ++(0,-0.8) -| (pix);
  \draw [arrow] (next) -- node[above] {nie} (pres);
  \draw [arrow] (pres) -- (save);
  \draw [arrow] (save) -- (stop);
\end{tikzpicture}
\caption{Schemat ogólny projekcji RTG}
\end{figure}

\section{Przykład użycia}

W pliku \texttt{dp\_testy.py} zdefiniowano scenariusz:
\begin{itemize}
  \item \texttt{create\_synthetic\_xray\_demo\_dicoms()},
  \item \texttt{build\_synthetic\_xray\_demo\_volumes()},
  \item \texttt{demo\_synthetic\_xray\_projection()}.
\end{itemize}

Przykład minimalny:

\begin{lstlisting}[style=python, caption={Minimalny scenariusz projekcji RTG}]
from dpVision import (
    XRayProjectionGeometry,
    XRayProjectionQualityProfile,
    XRayProjectionConfig,
    XRayPhysicsModel,
    DigitalRadiographyPresentationModel,
    VolumetricXRaySource,
    XRayScene,
    save_projection_png,
)

geometry = XRayProjectionGeometry.from_detector_pose(
    detector_center_ref=[42.2, 42.2, 180.0],
    detector_normal_ref=[0.0, 0.0, -1.0],
    detector_up_ref=[0.0, 1.0, 0.0],
    detector_shape_hw=[256, 256],
    detector_pixel_size_mm=0.4,
    step_mm=1.5,
    source_position_ref=[42.2, 42.2, -220.0],
)

physics = XRayPhysicsModel(output_mode="integral")
scene = XRayScene.from_sample_sources([
    VolumetricXRaySource(skull, global_transform=skull_transform),
    VolumetricXRaySource(jaw,   global_transform=jaw_transform),
])
config = XRayProjectionConfig(
    geometry=geometry,
    physics_model=physics,
    presentation_model=DigitalRadiographyPresentationModel(
        invert=False,
        gamma=0.7,
        contrast=1.2,
    ),
    quality_profile=XRayProjectionQualityProfile.normal(),
)

raw_image, stats = scene.project(config=config, return_stats=True)
display_image = config.apply_presentation(raw_image)

save_projection_png(display_image, "synthetic_xray.png",
                    invert=False, fixed_range=(0.0, 1.0))
\end{lstlisting}

\section{Zlozonosc obliczeniowa i koszt parametrow}

W biezacej implementacji koszt projekcji jest zdominowany przez ray-marching,
czyli przez liczbe promieni i liczbe probek pobieranych wzdluz kazdego promienia.
W pierwszym przyblizeniu mozna zapisac:
\begin{equation}
T \sim H \cdot W \cdot N_s \cdot N_v,
\end{equation}
gdzie:
\begin{itemize}
  \item \(H, W\) -- wysokosc i szerokosc siatki detektora,
  \item \(N_s\) -- liczba probek na promien,
  \item \(N_v\) -- liczba aktywnych zrodel danych, na przyklad wolumenow.
\end{itemize}

Jesli przecietna dlugosc odcinka promienia przebiegajacego przez scene wynosi \(L\),
a krok probkowania jest rowny \(\Delta s = \texttt{step\_mm}\), to:
\begin{equation}
N_s \approx \frac{L}{\Delta s},
\end{equation}
co prowadzi do praktycznej zaleznosci:
\begin{equation}
T \sim H \cdot W \cdot \frac{L}{\texttt{step\_mm}} \cdot N_v.
\end{equation}

\subsection{Co zwieksza czas liczenia}

\begin{enumerate}[label=\arabic*.]
  \item \textbf{Wieksza liczba pikseli detektora}
  
  Podwojenie rozdzielczosci w obu osiach daje w przyblizeniu czterokrotny wzrost
  liczby promieni:
  \begin{equation}
  (2H) \cdot (2W) = 4HW.
  \end{equation}

  \item \textbf{Mniejszy krok \texttt{step\_mm}}
  
  Zmniejszenie kroku z \(1.5\ \mathrm{mm}\) do \(0.75\ \mathrm{mm}\) daje
  w przyblizeniu dwukrotnie wiecej probek na promien:
  \begin{equation}
  \frac{L}{0.75} \approx 2 \cdot \frac{L}{1.5}.
  \end{equation}

  \item \textbf{Wieksza liczba wolumenow}
  
  W uogolnieniu czaszka + zuchwa sa drozsze od pojedynczej czaszki, poniewaz
  kazda probka promienia musi byc sprawdzona w wiecej niz jednym zrodle danych.
\end{enumerate}

\subsection{Co zmienia geometria, a co zmienia koszt}

Trzeba rozroznic dwa parametry o roznych znaczeniach:
\begin{itemize}
  \item \textbf{skala detektora}, wynikajaca z wektorow
        \(\mathbf{u}_{\mathrm{det}}\), \(\mathbf{v}_{\mathrm{det}}\) i z
        \texttt{detector\_shape\_hw},
  \item \textbf{krok integracji} \(\texttt{step\_mm}\), ktory nie zmienia
        pola widzenia, ale zmienia dokladnosc i czas ray-marchingu.
\end{itemize}

W aktualnym API dlugosci \(\|\mathbf{u}_{\mathrm{det}}\|\) i
\(\|\mathbf{v}_{\mathrm{det}}\|\) odpowiadaja rozmiarowi pojedynczego piksela
detektora w milimetrach. Z kolei fizyczny rozmiar detektora wynosi odpowiednio:
\begin{equation}
W_{\mathrm{det}} = W \cdot \|\mathbf{u}_{\mathrm{det}}\|,
\qquad
H_{\mathrm{det}} = H \cdot \|\mathbf{v}_{\mathrm{det}}\|.
\end{equation}

Zmiana rozmiaru detektora lub jego odleglosci od zrodla zmienia:
\begin{itemize}
  \item pole widzenia,
  \item powiekszenie perspektywiczne,
  \item to, jaki fragment sceny zmiesci sie na obrazie.
\end{itemize}

Nie musi jednak sama z siebie zmieniac czasu obliczen, jesli liczba pikseli
pozostaje taka sama. Czas rosnie wtedy, gdy zwiekszamy \texttt{shape\_hw},
a nie tylko fizyczny wymiar kliszy lub detektora.

\subsection{Praktyczne kompromisy}

\begin{itemize}
  \item \textbf{Do szybkiego debugowania geometrii:}
        mniejszy detektor, na przyklad \(128 \times 128\), oraz wiekszy
        \texttt{step\_mm}, na przyklad \(1.5\) -- \(2.0\ \mathrm{mm}\).
  \item \textbf{Do porownan scen lub animacji:}
        utrzymanie umiarkowanej rozdzielczosci detektora i stalego
        \texttt{step\_mm}, aby czasy kolejnych projekcji byly przewidywalne.
  \item \textbf{Do finalnego renderu:}
        wiekszy detektor, na przyklad \(256 \times 256\) lub \(512 \times 512\),
        oraz mniejszy krok, jesli potrzebna jest dokladniejsza integracja
        przez cienkie struktury kostne.
  \item \textbf{Do wielu obiektow ruchomych:}
        lepiej unikac kosztownego resamplingu posredniego i probkowac scene
        bezposrednio w przestrzeni swiata, bo dodatkowy resampling wprowadza
        i koszt, i kolejna warstwe interpolacji.
\end{itemize}

\section{Ograniczenia bieżącej implementacji}

\begin{itemize}
  \item Projekcja jest liczona na CPU i ma złożoność liniową względem liczby pikseli
        detektora, liczby próbek po promieniu i liczby źródeł.
  \item Model \(\mu(I_{\mathrm{CT}})\) jest prosty i nie reprezentuje jeszcze
        pełnej fizyki klinicznej.
  \item Detektor nie modeluje jeszcze odpowiedzi energetycznej, szumu, scatter ani MTF.
  \item Eksport DICOM używa na razie prostego \texttt{Secondary Capture}.
  \item Interfejs \texttt{XRaySampleSource} jest gotowy na siatki, ale bieżąca implementacja
        zawiera tylko adapter wolumetryczny.
\end{itemize}

\section{Kierunki rozbudowy}

\begin{enumerate}[label=\arabic*.]
  \item \textbf{Nowe źródła próbek:} siatki trójkątów, siatki materiałowe,
        źródła nieciągłe, maski segmentacyjne.
  \item \textbf{Lepsza fizyka:} osobne modele \(\mu(E)\), widmo lampy,
        beam hardening, scatter, szum kwantowy.
  \item \textbf{Akceleracja:} batching promieni, wektoryzacja, GPU.
  \item \textbf{Geometrie specjalne:} CBCT, panoramiczne, detektor zakrzywiony.
  \item \textbf{Integracja z GUI:} dialog ustawiania geometrii, podgląd i animacja
        scen czaszka + żuchwa.
\end{enumerate}

\section*{Podsumowanie}
\addcontentsline{toc}{section}{Podsumowanie}

Moduł \texttt{xrayProjection.py} nie jest tylko kolejną funkcją podglądu.
Stanowi oddzielny backend projekcji RTG, w którym:
\begin{itemize}
  \item geometria,
  \item źródła próbkowania,
  \item fizyka akwizycji,
  \item i prezentacja
\end{itemize}
są jawnie rozdzielone. Taki podział pozwala rozwijać zgodność z fizyką RTG
bez naruszania już działającego kodu dla wolumenów i bez przeciążania klasy
\texttt{Volumetric} odpowiedzialnościami niezwiązanymi z przechowywaniem danych.

\end{document}
